\manualchapter{Patches and troubleshooting}{sec:operation}
\subsection{External FM pair}
Patch a second Operator 2612 OUT into FM. Give both modules the same pitch and
gate, set the modulator's multiplier to 2, and raise the receiving module's
FM depth from zero. A short modulator envelope changes the tone at the start
of each note; a sustained modulator changes the whole note. Reduce either
FM depth or modulator volume if the pitch becomes hard to hear.

\subsection{When a control seems inactive}
Check GATE, TL and output volume first if there is no sound. A zero D1R
prevents the envelope reaching the second decay stage. AMS and FMS must be
nonzero to hear the internal LFO. The FM knob only has an effect with an
external signal. A patch to VOL is additive, so 0\,V does not mute a raised knob.

\subsection{Gates, polyphony and saved patches}
GATE opens at 2\,V and closes at or below 0.01\,V. RTRG uses the same
thresholds. Both are read every host sample: even a one-sample pulse is
recognized after a sampled low interval. A held-high RTRG causes one event.
Envelope parameter CV still updates every 16 samples; pitch updates every sample.

A rising RTRG restarts attack from the current chip attenuation. With GATE
held high, the key-off/key-on pair completes before the next audio sample.
Simultaneous gate and retrigger rises start once. With GATE low (including a
simultaneous gate fall), RTRG opens the key for one sample, then releases it.
Hold GATE high for continuing loops. Earlier builds sampled these inputs
every 16 samples; short pulses could be missed and retriggers could insert
16 samples of release. Existing patches may therefore have different note timing.

The widest input sets 1--16 polyphonic voices. Most inputs read the matching
channel directly; a mono gate does not automatically open every channel of
a polyphonic pitch cable. Supply matching gate and pitch channel counts.
Volume CV does repeat a mono voltage across active channels.

Parameters and the context option \textbf{Soft Reset Envelope Generator}
are saved in patches. The option is off by default and after reset. When
enabled, key-on preserves oscillator phase; when disabled, key-on resets it.
Both settings retain the current attenuation and the running envelope/LFO
clocks. Loop-boundary phase resets still occur with either setting; the
option does not guarantee click-free sound. Oscillator phase and envelope
progress are not saved.

Loop retains the existing simplified repeating SSG contour, not all hardware
SSG shapes. Slow attacks and zero decay rates keep their chip-style behavior:
a zero rate can hold a stage indefinitely, and a retrigger need not sound like
a fresh note from silence. These details are part of the sound, not reasons
to force the attack rate or clear the envelope level. Disconnected or removed
polyphonic lanes release their keys and rearm their input detectors.
